% Author: Grayson Orr
% Course: ID511001: Programming 2

\documentclass{article}
\author{}

\usepackage{graphicx}
\usepackage{wrapfig}
\usepackage{enumerate}
\usepackage{hyperref}
\usepackage[margin = 2.25cm]{geometry}
\usepackage[table]{xcolor}
\usepackage{soul}
\usepackage{fancyhdr}
\usepackage{fancyvrb}
\hypersetup{
  colorlinks = true,
  urlcolor = blue
}
\setlength\parindent{0pt}
\pagestyle{fancy}
\fancyhf{}
\rhead{College of Engineering, Construction \& Living Sciences\\Bachelor of Information Technology}
\lfoot{Project 1: Pong\\Version 1, Semester Two, 2022}
\rfoot{\thepage}
 
\begin{document}

\begin{figure}
    \centering
    \includegraphics[width=50mm]{../../resources/img/logo.png}
\end{figure}

\title{College of Engineering, Construction \& Living Sciences\\Bachelor of Information Technology\\ID511001: Programming 2\\Level 5, Credits 15\\\textbf{Project 1: Pong}}
\date{}
\maketitle

\section*{Assessment Overview}
In this assessment, you will design \& develop a GUI implementation of the classic arcade game \textbf{Pong}.

\section*{Learning Outcomes}
At the successful completion of this course, learners will be able to:
\begin{enumerate}
    \item Build interactive, event-driven GUI applications using pre-built components.
    \item Declare \& implement user-defined classes using encapsulation, inheritance \& polymorphism.
\end{enumerate}

\section*{Assessments}
\renewcommand{\arraystretch}{1.5}
\begin{tabular}{|c|c|c|c|}
	\hline
	\textbf{Assessment}                                 & \textbf{Weighting} & \textbf{Due Date}            & \textbf{Learning Outcomes} \\ \hline
	\small Project 1: Pong & \small 25\%        & \small 17-10-2022 (Monday at 4.59 PM)   & \small 1 \& 2                   \\ \hline
	\small Project 2: Space Invaders & \small 35\%        & \small 04-11-2022 (Friday at 4.59 PM)   & \small 1 \& 2                   \\ \hline
	\small Theory Examination                        & \small 30\%        & \small 18-11-2022 (Friday at 4.59 PM)  & \small 1 \& 2                   \\ \hline
	\small Classroom Tasks                       & \small 10\%        & \small 21-10-2022 (Friday at 4.59 PM)  & \small 1 \& 2                   \\ \hline
\end{tabular}

\section*{Conditions of Assessment}
You will complete this assessment during your learner-managed time. However, there will be time during class to discuss the requirements \& your progress on this assessment. This assessment will need to be completed by \textbf{Monday, 17 October 2022} at \textbf{4.59 PM}.

\section*{Pass Criteria}
This assessment is criterion-referenced (CRA) with a cumulative pass mark of \textbf{50\%} over all assessments in \textbf{ID511001: Programming 2}.

\section*{Authenticity}
All parts of your submitted assessment \textbf{must} be completely your work. If you use code snippets from \textbf{GitHub}, \textbf{StackOverflow} or other online resources, you \textbf{must} reference it appropriately using \textbf{APA 7th edition}. Provide your references in the \textbf{README.md} file in your repository. Failure to do this will result in a mark of \textbf{zero} for this assessment.

\section*{Policy on Submissions, Extensions, Resubmissions \& Resits}
The school's process concerning submissions, extensions, resubmissions \& resits complies with \textbf{Otago Polytechnic} policies. Learners can view policies on the \textbf{Otago Polytechnic} website located at \href{https://www.op.ac.nz/about-us/governance-and-management/policies}{https://www.op.ac.nz/about-us/governance-and-management/policies}.

\section*{Submission}
You \textbf{must} submit all project files via \textbf{GitHub Classroom}. Here is the URL to the repository you will use for your submission – \href{https://classroom.github.com/a/Kjpy9SCn}{https://classroom.github.com/a/Kjpy9SCn}.  Create a \textbf{.gitignore} \& add the ignored files in this resource - \href{https://raw.githubusercontent.com/github/gitignore/main/VisualStudio.gitignore}{https://raw.githubusercontent.com/github/gitignore/main/VisualStudio.gitignore}. The latest project files in the \textbf{master} or \textbf{main} branch will be used to mark against the \textbf{Functionality} criterion. Please test before you submit. Partial marks \textbf{will not} be given for incomplete functionality. Late submissions will incur a \textbf{10\% penalty per day}, rolling over at \textbf{5:00 PM}.

\section*{Extensions}
Familiarise yourself with the assessment due date. Contact the course lecturer before the due date if you need an extension. If you require more than a week's extension, you will need to provide a medical certificate or support letter from your manager.

\section*{Resubmissions}
Learners may be requested to resubmit an assessment following a rework of part/s of the original assessment. Resubmissions are to be completed within a negotiable short time frame \& usually \textbf{must} be completed within the timing of the course to which the assessment relates. Resubmissions will be available to learners who have made a genuine attempt at the first assessment opportunity \& achieved a \textbf{D grade (40-49\%)}. The maximum grade awarded for resubmission will be \textbf{C-}.

\section*{Resits}
Resits \& reassessments \textbf{are not} applicable in \textbf{ID511001: Programming 2}.

\section*{Instructions}
You will need to submit a project \& documentation that meet the following requirements:\\

\textbf{Note:} Independent research requirements are \hl{highlighted yellow}.

\subsection*{Functionality - Learning Outcomes 1 \& 2 (40\%)}
\begin{itemize}
    \item Project \textbf{must} open with code or file structure modification in \textbf{Visual Studio}.
    \item A game of \textbf{Pong} \textbf{must} be driven by one \textbf{Timer}.
    \item A ball \& two paddles \textbf{must} be created using the \textbf{Graphics} class.
    \item The ball \textbf{must} bounce (collide) off the top \& bottom of the screen, \& paddles.
    \item The paddles \textbf{must} move vertically but not exceed the top \& bottom of the screen.
    \item The user controls the left paddle via the up \& down keys. The computer controls the right paddle. It is acceptable for the right paddle to follow the ball's position. However, other solutions are encouraged.
    \item A scoring system. When the ball collides with the left \& right-hand side of the screen, one point is given to either the user or computer. The game is over when either score is 10. Appropriate feedback \textbf{must} be displayed to the user, i.e., \textbf{"You win!"} or \textbf{"You lose!"}.
    \item Double buffering to prevent the ball \& paddles from flickering.
    \item \hl{\textbf{Independent Research:}}
    \begin{itemize}
        \item A highscore system. When the game is over, the user \& computer's scores are saved, i.e., written to a text file. Read the user's scores from the text file and display the last five to the user.
        \item Play a sound when:
        \begin{itemize}
            \item The ball bounces off the paddle, \& top \& bottom of the screen.
            \item The user wins.
            \item The user loses.
        \end{itemize}
        \textbf{Note:} These sounds \textbf{must} be unique.
        \item An ability to play a new game, restart a current game \& pause a current game.
        \item Randomise the colour of the ball \& paddles.
    \end{itemize}
\end{itemize}

\subsection*{Code Elegance - Learning Outcomes 1 \& 2  (45\%)}
\begin{itemize}
    \item Adhere to the four principles of \textbf{OO}, i.e., encapsulation, abstraction, inheritance \& polymorphism.
    \item Use of intermediate variables, constants \& enumerations.
    \item Idiomatic use of control flow, data structures \& in-built functions.
    \item Efficient algorithmic approach.
    \item Sufficient modularity.
    \item Each method \& class \textbf{must} have a header comment located immediately before its declaration.
    \item In-line comments where required. 
    \item Project files, i.e., \textbf{.cs} files are fomatted. 
    \item No dead or unused code.
\end{itemize}

\subsection*{Documentation \& Git Usage - Learning Outcomes 1 \& 2 (15\%)}
\begin{itemize}
    \item Provide the following in your repository \textbf{README.md} file:
    \begin{itemize}
        \item Your project's UML diagram.
        \item References to used code snippets from \textbf{GitHub}, \textbf{StackOverflow} or other online resources.
        \item Known bugs if applicable.
    \end{itemize}
    \item Commit at least \textbf{20} times per week.
    \item Commit messages \textbf{must} reflect the context of each functional requirement change \& formatted using the naming conventions dicussed in \textbf{Week 1}.
\end{itemize}

\subsection*{Additional Information}
\begin{itemize}
    \item \textbf{Do not} rewrite your \textbf{Git} history. It is important that the course lecturer can see how you worked on your assessment over time.
    \item When the user presses a key, i.e., up or down, a \textbf{KeyDown} event is generated. For the \textbf{Form1's} \textbf{KeyDown} event, the method signature is:
    \begin{Verbatim}[tabsize=2]
		private void Form1_KeyDown(object sender, KeyEventArgs e) {}
	\end{Verbatim}
	The argument you will be interested in is \textbf{KeyEventArgs e} which is the value of the pressed key. The arrow key values are \textbf{Keys.Left}, \textbf{Keys.Right}, \textbf{Keys.Up} \& \textbf{Keys.Down}. In the \textbf{Form1\_KeyDown} method, you can use a \textbf{switch} statement. For example:
	\begin{Verbatim}[tabsize=2]
		switch (e.KeyCode) 
		{
			case Keys.Left:
				// Do something
				break;
			case Keys.Right:
				// Do something
				break;
			case Keys.Up:
				// Do something
				break;
			case Keys.Down:
				// Do something
				break;
			default:
				// Do something
				break;
		}
	\end{Verbatim}
	\textbf{Note:} The \textbf{Form1's} \textbf{KeyPreview} event \textbf{must} be set to \textbf{True}. Otherwise, \textbf{Form1} will not respond to the \textbf{KeyDown} event.
\end{itemize}

\end{document}